\section{Características generales de las células eucariotas}

Las células eucariotas son organismos (unicelulares o pluricelulares) que poseen núcleo definido y organelos membranosos especializados. Presentan un tamaño mayor que las procariotas (10-100 µm).

\section{Tipos de células eucariotas}

\subsection{Célula animal}
\begin{itemize}
    \item Sin pared celular
    \item Sin cloroplastos
    \item Presenta centriolos
    \item Vacuolas pequeñas y numerosas
    \item Membrana plasmática como límite externo
\end{itemize}

\subsection{Célula vegetal}
\begin{itemize}
    \item Pared celular de celulosa
    \item Cloroplastos para fotosíntesis
    \item Vacuola central grande
    \item Sin centriolos
    \item Plasmodesmos para comunicación
\end{itemize}

\subsection{Célula fúngica}
\begin{itemize}
    \item Pared celular de quitina
    \item Vacuolas prominentes
    \item Hifas (estructura filamentosa)
    \item Heterótrofos por absorción
\end{itemize}

\subsection{Célula protista}
\begin{itemize}
    \item Organismos eucariotas unicelulares
    \item Diversidad morfológica
    \item Pueden ser autótrofos o heterótrofos
    \item Ejemplos: Amoeba, Paramecium, Euglena
\end{itemize}

\section{Organelos celulares y sus funciones}

\begin{table}[h]
\centering
\begin{tabular}{|l|p{6cm}|}
\hline
\textbf{Organelo} & \textbf{Función principal} \\
\hline
Núcleo & Almacenamiento y transmisión de información genética \\
\hline
Mitocondrias & Respiración celular y producción de ATP \\
\hline
Retículo endoplásmico & Síntesis y transporte de proteínas y lípidos \\
\hline
Aparato de Golgi & Procesamiento y empaquetamiento de proteínas \\
\hline
Lisosomas & Digestión celular y reciclaje de componentes \\
\hline
Peroxisomas & Oxidación de ácidos grasos y detoxificación \\
\hline
Cloroplastos & Fotosíntesis (en células vegetales) \\
\hline
Vacuolas & Almacenamiento y mantenimiento de turgencia \\
\hline
Ribosomas & Síntesis de proteínas \\
\hline
\end{tabular}
\caption{Organelos de la célula eucariota y sus funciones}
\end{table}

\section{Diferencias entre célula animal y vegetal}

\begin{table}[h]
\centering
\begin{tabular}{|l|c|c|}
\hline
\textbf{Característica} & \textbf{Animal} & \textbf{Vegetal} \\
\hline
Pared celular & No & Sí (celulosa) \\
\hline
Cloroplastos & No & Sí \\
\hline
Vacuola central & No & Sí \\
\hline
Centriolos & Sí & No \\
\hline
Plasmodesmos & No & Sí \\
\hline
Plástidos & No & Sí \\
\hline
\end{tabular}
\caption{Comparación entre célula animal y vegetal}
\end{table}